% GENERATED FILE. Edit canonical lessons, then run npm run generate:books.
% canonical-source-hash: fnv1a64:7f34423fc9036ea3
% canonical-lessons: MW-C201-kot, MW-C201-mohallo, MW-C201-thano, MW-C201-saray, MW-C201-ghusno

\chapter{Around the Old Town}
\label{ch:mw-around-the-old-town}
\hlchaptermodality{201}

\begin{chapteropening}
\textbf{By the end of this chapter:} \emph{I can ask where a place in town is.}

Complete the last lesson of chapter 201: I can ask where a place in town is.
\end{chapteropening}

\section[\texorpdfstring{\mw{कोट} (koṭ)}{कोट (koṭ)}]{\mw{कोट} (koṭ) — a fort}

\label{lesson:MW-C201-kot}



\begin{quote}
Before the new one: say the Marwadi for a quarter, then the Marwadi for a litre.
\end{quote}

\subsection*{You'll want to know: \mw{कोट}}
\textbf{\mw{कोट}} — \emph{koṭ} — \textquotedblleft{}a fort\textquotedblright{}.

\begin{quote}
\mw{कोट} \mw{कठै} \mw{है}?
\end{quote}

\begin{quote}
\emph{koṭ kaṭhai hai?}
\end{quote}

\begin{quote}
Where is the fort?
\end{quote}

A masculine noun.

\begin{cousinweb}
\textbf{\mw{कोट}} (\emph{koṭ}) is inherited from Sanskrit \emph{kōṭṭa}, \textquotedblleft{}fort\textquotedblright{}; CDIAL lists Marwari \emph{koṭ}, \textquotedblleft{}fort\textquotedblright{} (CDIAL 3500).
\end{cousinweb}

\begin{grammarlens}[title={What you've built}]
One more word for this chapter.
\end{grammarlens}

\subsection*{Your turn}
\begin{itemize}
\raggedright
  \item \emph{Say it:} \emph{koṭ}
  \item \emph{Say it:} \emph{koṭ}, once more
  \item \emph{Say it:} \emph{koṭ kaṭhai hai?} — Where is the fort?
  \item \emph{Recall:} say the Marwadi for a quarter, then the Marwadi for a litre, then say \emph{koṭ} again
\end{itemize}

\begin{tcolorbox}[breakable,colback=teal!4,colframe=teal!35!black,title={Before you move on}]
What does \mw{कोट} mean? (\textquotedblleft{}A fort\textquotedblright{}.) Say it once more.
\end{tcolorbox}
\section[\texorpdfstring{\mw{मोहल्लो} (mohallo)}{मोहल्लो (mohallo)}]{\mw{मोहल्लो} (mohallo) — a neighbourhood}

\label{lesson:MW-C201-mohallo}



\begin{quote}
Before the new one: say the Marwadi for a litre, then the Marwadi for a fort.
\end{quote}

\subsection*{You'll want to know: \mw{मोहल्लो}}
\textbf{\mw{मोहल्लो}} — \emph{mohallo} — \textquotedblleft{}a neighbourhood\textquotedblright{}.

\begin{quote}
\mw{म्हारो} \mw{मोहल्लो} \mw{कोट} \mw{कनै} \mw{है।}
\end{quote}

\begin{quote}
\emph{mhāro mohallo koṭ kanai hai.}
\end{quote}

\begin{quote}
My neighbourhood is near the fort.
\end{quote}

A masculine noun.

\begin{cousinweb}
\textbf{\mw{मोहल्लो}} (\emph{mohallo}) is usually traced, through Persian, to Arabic \emph{maḥalla}, \textquotedblleft{}a quarter of a town\textquotedblright{}.
\end{cousinweb}

\begin{grammarlens}[title={What you've built}]
One more word for this chapter.
\end{grammarlens}

\subsection*{Your turn}
\begin{itemize}
\raggedright
  \item \emph{Say it:} \emph{mohallo}
  \item \emph{Say it:} \emph{mohallo}, once more
  \item \emph{Say it:} \emph{mhāro mohallo koṭ kanai hai.} — My neighbourhood is near the fort.
  \item \emph{Recall:} say the Marwadi for a litre, then the Marwadi for a fort, then say \emph{mohallo} again
\end{itemize}

\begin{tcolorbox}[breakable,colback=teal!4,colframe=teal!35!black,title={Before you move on}]
What does \mw{मोहल्लो} mean? (\textquotedblleft{}A neighbourhood\textquotedblright{}.) Say it once more.
\end{tcolorbox}
\section[\texorpdfstring{\mw{थाणो} (thāṇo)}{थाणो (thāṇo)}]{\mw{थाणो} (thāṇo) — a police station}

\label{lesson:MW-C201-thano}



\begin{quote}
Before the new one: say the Marwadi for a fort, then the Marwadi for a neighbourhood.
\end{quote}

\subsection*{You'll want to know: \mw{थाणो}}
\textbf{\mw{थाणो}} — \emph{thāṇo} — \textquotedblleft{}a police station\textquotedblright{}.

\begin{quote}
\mw{थाणो} \mw{चौक} \mw{कनै} \mw{है।}
\end{quote}

\begin{quote}
\emph{thāṇo chauk kanai hai.}
\end{quote}

\begin{quote}
The police station is near the square.
\end{quote}

A masculine noun.

\begin{cousinweb}
\textbf{\mw{थाणो}} (\emph{thāṇo}) is inherited from Sanskrit \emph{sthāna}, \textquotedblleft{}a place\textquotedblright{}. CDIAL gives no Marwari form in this entry, but lists Hindi \emph{thānā} and Gujarati \emph{thāṇũ}, \textquotedblleft{}police station\textquotedblright{} (CDIAL 13753).
\end{cousinweb}

\begin{grammarlens}[title={What you've built}]
One more word for this chapter.
\end{grammarlens}

\subsection*{Your turn}
\begin{itemize}
\raggedright
  \item \emph{Say it:} \emph{thāṇo}
  \item \emph{Say it:} \emph{thāṇo}, once more
  \item \emph{Say it:} \emph{thāṇo chauk kanai hai.} — The police station is near the square.
  \item \emph{Recall:} say the Marwadi for a fort, then the Marwadi for a neighbourhood, then say \emph{thāṇo} again
\end{itemize}

\begin{tcolorbox}[breakable,colback=teal!4,colframe=teal!35!black,title={Before you move on}]
What does \mw{थाणो} mean? (\textquotedblleft{}A police station\textquotedblright{}.) Say it once more.
\end{tcolorbox}
\section[\texorpdfstring{\mw{सराय} (sarāy)}{सराय (sarāy)}]{\mw{सराय} (sarāy) — an inn}

\label{lesson:MW-C201-saray}



\begin{quote}
Before the new one: say the Marwadi for a neighbourhood, then the Marwadi for a police station.
\end{quote}

\subsection*{You'll want to know: \mw{सराय}}
\textbf{\mw{सराय}} — \emph{sarāy} — \textquotedblleft{}an inn\textquotedblright{}.

\begin{quote}
\mw{सराय} \mw{टेसण} \mw{सामी} \mw{है।}
\end{quote}

\begin{quote}
\emph{sarāy ṭesaṇ sāmī hai.}
\end{quote}

\begin{quote}
The inn is opposite the station.
\end{quote}

A feminine noun.

\begin{cousinweb}
\textbf{\mw{सराय}} (\emph{sarāy}) is usually traced to Persian \emph{sarāy}, \textquotedblleft{}a palace; an inn\textquotedblright{}.
\end{cousinweb}

\begin{grammarlens}[title={What you've built}]
One more word for this chapter.
\end{grammarlens}

\subsection*{Your turn}
\begin{itemize}
\raggedright
  \item \emph{Say it:} \emph{sarāy}
  \item \emph{Say it:} \emph{sarāy}, once more
  \item \emph{Say it:} \emph{sarāy ṭesaṇ sāmī hai.} — The inn is opposite the station.
  \item \emph{Recall:} say the Marwadi for a neighbourhood, then the Marwadi for a police station, then say \emph{sarāy} again
\end{itemize}

\begin{tcolorbox}[breakable,colback=teal!4,colframe=teal!35!black,title={Before you move on}]
What does \mw{सराय} mean? (\textquotedblleft{}An inn\textquotedblright{}.) Say it once more.
\end{tcolorbox}
\section[\texorpdfstring{\mw{घुसणो} (ghusṇo)}{घुसणो (ghusṇo)}]{\mw{घुसणो} (ghusṇo) — to get in; to enter}

\label{lesson:MW-C201-ghusno}



\begin{quote}
Before the new one: say the Marwadi for a police station, then the Marwadi for an inn.
\end{quote}

\subsection*{You'll want to know: \mw{घुसणो}}
\textbf{\mw{घुसणो}} — \emph{ghusṇo} — \textquotedblleft{}to get in; to enter\textquotedblright{}.

\begin{quote}
\mw{बकरी} \mw{खेत} \mw{मांय} \mw{घुसै} \mw{है।}
\end{quote}

\begin{quote}
\emph{bakrī khet mā̃y ghusai hai.}
\end{quote}

\begin{quote}
The goat keeps getting into the field.
\end{quote}

Past (he …): \mw{वो} \mw{मांय} \mw{घुस्यो}, \emph{vo mā̃y ghusyo}.

\begin{cousinweb}
\textbf{\mw{घुसणो}} (\emph{ghusṇo}) is inherited from a reconstructed \emph{*ghuss}, \textquotedblleft{}thrust in\textquotedblright{}; CDIAL lists Marwari \emph{ghusṇo}, \textquotedblleft{}to enter\textquotedblright{} (CDIAL 4492).
\end{cousinweb}

\begin{grammarlens}[title={What you've built}]
One more word for this chapter.
\end{grammarlens}

\subsection*{Your turn}
\begin{itemize}
\raggedright
  \item \emph{Say it:} \emph{ghusṇo}
  \item \emph{Say it:} \emph{ghusṇo}, once more
  \item \emph{Say it:} \emph{bakrī khet mā̃y ghusai hai.} — The goat keeps getting into the field.
  \item \emph{Recall:} say the Marwadi for a police station, then the Marwadi for an inn, then say \emph{ghusṇo} again
\end{itemize}

\begin{tcolorbox}[breakable,colback=teal!4,colframe=teal!35!black,title={Before you move on}]
What does \mw{घुसणो} mean? (\textquotedblleft{}To get in; to enter\textquotedblright{}.) Say it once more.
\end{tcolorbox}
